\documentclass[12pt,a4paper]{article}

\usepackage[utf8]{inputenc}
\usepackage[T1]{fontenc}
\usepackage[brazil]{babel}
\usepackage{amsmath,amssymb}
\usepackage{algorithm}
\usepackage{algpseudocode}
\usepackage{geometry}

\geometry{margin=2.5cm}

\title{Algoritmo Espectral para Verificação de Bipartição de Grafos}
\author{Prof. Me. Luis Vinicius Costa Silva}
\date{}

\begin{document}

\maketitle

\section{Problema}

Considere um grafo simples e não direcionado

$$
G=(V,E),
$$

em que cada vértice representa um elemento e cada aresta representa
um conflito entre dois elementos.

O objetivo é determinar se os vértices podem ser divididos em dois
grupos de modo que nenhum par de vértices conflitantes pertença ao
mesmo grupo.

Esse problema é equivalente a determinar se o grafo é
\emph{bipartido}.

\section{Representação por matriz de adjacência}

Seja $A\in\mathbb{R}^{n\times n}$ a matriz de adjacência de $G$,
definida por

$$
A_{ij} =
\begin{cases}
1, & \text{se } (i,j)\in E,\\
0, & \text{caso contrário}.
\end{cases}
$$

Como o grafo é não direcionado, temos

$$
A=A^T.
$$

Consequentemente, todos os autovalores de $A$ são reais.

\section{Caracterização espectral}

Um resultado fundamental da teoria espectral de grafos estabelece que
um grafo não direcionado é bipartido se, e somente se, o espectro de
sua matriz de adjacência é simétrico em relação à origem.

Se

$$
\lambda_1,\lambda_2,\ldots,\lambda_n
$$

são os autovalores de $A$, então $G$ é bipartido se, e somente se, para
cada autovalor $\lambda$, também existe o autovalor $-\lambda$, com a
mesma multiplicidade.

Assim,

$$
G\text{ é bipartido}
\iff
\operatorname{spec}(A)
=
-\operatorname{spec}(A).
$$

Após ordenar os autovalores de forma crescente,

$$
\lambda_1\leq\lambda_2\leq\cdots\leq\lambda_n,
$$

a condição pode ser verificada por

$$
\lambda_i=-\lambda_{n+1-i},
\qquad
i=1,\ldots,n.
$$

Em uma implementação numérica, devido aos erros de ponto flutuante,
utiliza-se uma tolerância $\varepsilon>0$:

$$
\left|\lambda_i+\lambda_{n+1-i}\right|<\varepsilon.
$$

\section{Justificativa}

Suponha que $G$ seja bipartido, com partições $V_1$ e $V_2$.
Reordenando os vértices de modo que os vértices de $V_1$ apareçam
primeiro e os de $V_2$ depois, sua matriz de adjacência possui a forma

$$
A=
\begin{pmatrix}
0 & B\\
B^T & 0
\end{pmatrix}.
$$

Considere um autovetor

$$
\begin{pmatrix}
x\\y
\end{pmatrix}
$$

associado ao autovalor $\lambda$. Então

$$
A
\begin{pmatrix}
x\\y
\end{pmatrix}
=
\lambda
\begin{pmatrix}
x\\y
\end{pmatrix}.
$$

Isso implica

$$
By=\lambda x
$$

e

$$
B^Tx=\lambda y.
$$

Consequentemente,

$$
A
\begin{pmatrix}
x\\-y
\end{pmatrix}
=
-\lambda
\begin{pmatrix}
x\\-y
\end{pmatrix}.
$$

Portanto, sempre que $\lambda$ é autovalor, $-\lambda$ também é
autovalor. O espectro é, assim, simétrico em relação a zero.

O resultado inverso também é válido: se o espectro da matriz de
adjacência é simétrico em relação a zero, então o grafo é bipartido.

\section{Algoritmo}

A abordagem espectral consiste nas seguintes etapas:

\begin{enumerate}
\item construir a matriz de adjacência $A$;
\item calcular todos os autovalores de $A$;
\item ordenar os autovalores;
\item comparar cada autovalor com o seu correspondente oposto;
\item se todos os pares forem compatíveis, declarar que a divisão
em dois grupos é possível;
\item caso contrário, declarar que a divisão é impossível.
\end{enumerate}

\begin{algorithm}
\caption{Verificação espectral de bipartição}
\begin{algorithmic}[1]

\Require Grafo $G=(V,E)$ e tolerância $\varepsilon$
\Ensure Possibilidade ou impossibilidade de divisão em dois grupos

\State Construir a matriz de adjacência $A$
\State Calcular os autovalores

$$
\lambda_1,\ldots,\lambda_n
$$

de $A$
\State Ordenar os autovalores de forma crescente

\For{$i=1$ até $n$}
\If{$|\lambda_i+\lambda_{n+1-i}|>\varepsilon$}
\State \Return \textbf{DIVISÃO IMPOSSÍVEL}
\EndIf
\EndFor

\State \Return \textbf{DIVISÃO POSSÍVEL}

\end{algorithmic}
\end{algorithm}

\section{Exemplo}

Considere o caminho com quatro vértices:

$$
0-1-2-3.
$$

Sua matriz de adjacência é

$$
A=
\begin{pmatrix}
0&1&0&0\\
1&0&1&0\\
0&1&0&1\\
0&0&1&0
\end{pmatrix}.
$$

Os autovalores são aproximadamente

$$
-1.618,\quad
-0.618,\quad
0.618,\quad
1.618.
$$

Observa-se que

$$
-1.618\approx -0.618? 
$$

Mais precisamente, os pares simétricos são

$$
-1.618\leftrightarrow1.618
$$

e

$$
-0.618\leftrightarrow0.618.
$$

Logo, o espectro é simétrico em relação a zero e o grafo é bipartido.

Uma possível divisão é

$$
V_1=\{0,2\},
\qquad
V_2=\{1,3\}.
$$

Portanto,

$$
\boxed{\text{DIVISÃO POSSÍVEL}}
$$

\section{Implementação numérica}

Como a matriz de adjacência de um grafo não direcionado é simétrica,
pode-se utilizar o método de Jacobi para calcular seus autovalores.

O método aplica sucessivas rotações ortogonais à matriz até que os
elementos fora da diagonal sejam suficientemente pequenos. Ao final,
os elementos da diagonal fornecem aproximações dos autovalores.

Na implementação em C, pode-se utilizar uma condição de convergência
como

$$
\max_{i\neq j}|A_{ij}|<\varepsilon.
$$

Os autovalores obtidos são então ordenados e submetidos ao teste de
simetria descrito anteriormente.

\section{Complexidade}

Para uma matriz de adjacência de dimensão $n$, o armazenamento requer

$$
O(n^2)
$$

posições de memória.

O cálculo dos autovalores por um método iterativo como Jacobi possui
custo computacional superior ao da solução baseada em busca em largura.
Em contraste, a verificação de bipartição por BFS possui complexidade

$$
O(|V|+|E|)
$$

quando o grafo é representado por lista de adjacência.

Assim, a abordagem espectral não é necessariamente a mais eficiente
para decidir bipartição. Sua principal vantagem, neste contexto, é
demonstrar a relação entre a estrutura do grafo e o espectro de sua
matriz de adjacência.

\section{Conclusão}

O problema de dividir os vértices em dois grupos sem conflitos é
equivalente ao problema de verificar se o grafo é bipartido.

A abordagem espectral resolve essa questão por meio de uma propriedade
dos autovalores da matriz de adjacência:

$$
\boxed{
G\text{ é bipartido}
\iff
\operatorname{spec}(A)
=
-\operatorname{spec}(A)
}
$$

Dessa forma, a solução transforma um problema combinatório de
coloração em um problema de álgebra linear envolvendo os autovalores
da matriz de adjacência.

\end{document}
